% - [x] wymyśleć jakieś sukcesy z metrykami
% - [x] appreciated for / praised for
% - [x] Soft skille w przykładach
% - [?] top 3 rzeczy które po zatrudnieniu mogę wnieść do zespołu
% - [ ] We -> I w CSie
% - [ ] Wspomnieć we wstępie o sposobie myślenia, komunikacji z komputerem itd, myślenie schematami
% - [ ] Wspomnieć o Quant Strats Scholarship w uczelni i CS

\documentclass[a4paper, 12pt]{article} % a4 = 210mm * 297mm

\usepackage[T1]{fontenc}
\usepackage[a4paper, left=12pt, top=36pt, right=12pt, bottom=12pt]{geometry}
\usepackage[hidelinks]{hyperref}
\usepackage%[stretch = 25, shrink = 25, tracking=true, letterspace=30]%
{microtype}
\usepackage{xcolor}
\usepackage[fixed]{fontawesome5}

\definecolor{lightcolor}{HTML}{e9e4f8}
\definecolor{darkcolor}{HTML}{625892}
\newenvironment{items}[1]{\centering\textbf{\color{darkcolor}\large#1} \begin{itemize} \setlength{\itemsep}{0pt}}{\end{itemize}}
% All sans serif
\renewcommand{\familydefault}{\sfdefault}

% See overfull boxes better
\overfullrule10pt

% \setlength{\topskip}{0pt}
% \setlength{\parindent}{0pt}
\setlength{\fboxsep}{0pt}
% \setlength{\itemsep}{0pt}
\pagestyle{empty}

% Define minipage sizes
\newlength\pagesize
\setlength\pagesize{\textwidth}
\newlength{\leftwidth}
\setlength{\leftwidth}{0.3\pagesize}
\newlength{\rightwidth}
\setlength{\rightwidth}{\pagesize}
\addtolength{\rightwidth}{-\leftwidth}
\newlength\minispace
\setlength\minispace{12pt}
\addtolength{\leftwidth}{-0.5\minispace}
\addtolength{\rightwidth}{-0.5\minispace}

\begin{document}
\textbf{\hfill\Huge\color{darkcolor}Borys Kopeć\hfill}
\vspace{24pt}

\noindent\colorbox{lightcolor}{%
\begin{minipage}[t][][s]{\leftwidth}
  \vspace{1em}
  \begin{description}
    \item \faPhoneSquare \texttt{+}48 987 654 321
    \item \faEnvelopeSquare boryskopec00@gmail.com
    \item \faLinkedin \href{https://linkedin.com/in/boryskopec}{\texttt{/in/boryskopec}}
    \item \faGithubSquare \href{https://github.com/B0ryskart0n}{\texttt{/B0ryskart0n}}
  \end{description}
  \begin{items}{Software development}
    \item Agile manifesto
    \item Object-oriented\newline programming
    \item Functional programming
    \item Continous integration\newline Continous delivery
  \end{items}
  \begin{items}{Technologies}
    \item Rust
    \item F\texttt{\#}
    \item Visual Studio (Code)
    \item Microsoft Excel
    \item Git
    \item Linux
  \end{items}
  \begin{items}{Soft skills}
    \item Excellent memory
    \item Knowledge sharing
    \item Patience
    \item Public speaking
    \item Critical thinking
    \item Prioritzation and\newline working under pressure
  \end{items}
  \begin{items}{Privately}
    \item Developing my own\newline roguelike
    \item Board-gaming with\newline friends and family
  \end{items}
  \vspace{1em}
\end{minipage}}%
\hskip\minispace%
\begin{minipage}[t][][b]{\rightwidth}%
  \setlength{\parskip}{4pt}
  I enjoy applying maths, and programming is my favorite way of doing that.
  I also love games, both board and computer ones, and my free time revolves around those.

  I'm a very curious and critical person.
  I ask a lot of questions and I always try to get to the bottom of things.

  My favourite part of pretty much any work is supporting people and sharing knowledge.
  I'm always eager to help.

  \vspace{-3\parskip}
  \section*{\centering\color{darkcolor}Professional experience}
  {\textbf{Model Validator {\sf at} Commerzbank}\hfill\textcolor{gray}{\sl VIII 2024 -- present}}
  \vspace{\parskip}
  
  Validating Counterparty Credit Risk models.
  Performing analyses using \textbf{R}, \textbf{Excel}, and \textbf{SQL}, writing reports in \textbf{\LaTeX}.
  Working with model developers on one side, and regulators on the other side.

  So far, I've performed several regular validations, such as: IMM coverage, quarterly ISDA SIMM backtesting.
  I've had solid success at coordinating validation results with model owners, which required a good balance of regulation-strictness and pragmatism.

  Constantly working on streamlining and optimising validation work.
  
  \vspace{3\parskip}
  {\textbf{Quantitative Developer {\sf at} Credit Suisse}\hfill\textcolor{gray}{\sl VI 2022 -- VII 2024}}
  \vspace{\parskip}
  
  Worked in a team responsible for developing (\textbf{F\texttt{\#}} \& \textbf{C\texttt{\#}}) and supporting front-office applications and services.
  The main tool my team was responsible for was a desktop application for portfolio pricing and scenario risk analysis.
  
  My colleagues called me the "Master of Support".
  My number of quarterly handled support cases was around 15, while the median was around 3.

  During the turbulent times of UBS acquisition we experienced a shift in responsibilities towards the portfolio migration.
  We added support for several new trade types in the span of weeks.

  \vspace{-3\parskip}
  \section*{\centering\color{darkcolor} Education}
  {\textbf{Bachelor of Science in Applied Mathematics {\sf at} Wrocław University of Science and Technology}\hfill\textcolor{gray}{\sl X 2019 -- II 2023}}
  \vspace{\parskip}

  Learned on various mathematical topics, concurrently implementing most of the learnt concepts.
  I've been writing apps, analyses, tests, and simulations in \textbf{Python}, \textbf{R}, and \textbf{Julia}.
  I was primarily interested in Mathematical Modelling and Probability Theory.
  I've concluded the degree with my thesis "Estimating option prices with the Heston model".

  Repeated WUST Principal's Scholarship recipient. Year representative.
  Proudly selected for QuantStrats Scholarship Programme organised by Credit Suisse at which I got offered a Quantitative Developer position (described above).
\end{minipage}
\vspace{1em}

\noindent\textit{\footnotesize This document has been crafted by me, with no AI usage, in \LaTeX.}

\let\thefootnote\relax
\let\footnoterule\relax
\footnote{I hereby give consent for my personal data included in this document to be used for the purposes of the recruitment process.}
\end{document}
